\section{Skills Developed}

This project required the integration of computer vision (YOLOv8n detection, ROI geometry, temporal persistence), video coding (FFmpeg, HEVC, GOP/bitrate control, PSNR/SSIM/VMAF metrology), full-stack development (Python edge services, Node.js/Express servers, Next.js dashboards), experimental methodology (frozen protocols, provenance recording, group-aware splitting, pre-registered forensic expectations, BD-rate withholding criteria), and academic writing under strict evidence discipline, including the reporting of negative results.

\section{Challenges Faced and Overcome}

The principal challenges were methodological rather than purely technical: concentrating all calibration positives in one video forced an honest degeneracy analysis instead of a headline metric; the foreground/background split's collapse on moving-camera content forced a negative-result narrative instead of a success story; and insufficient rate--distortion overlap forced BD-rate to be withheld rather than resurrected from older logs. Each was resolved by weakening the claim to match the evidence, which is the central professional lesson of the work.

\section{Applications of Learning}

The evidence-discipline practices adopted here---provenance files, evidence matrices, claim audits, number-consistency checks, and explicit NOT AVAILABLE markings---transfer directly to any future engineering or research role in which measurements inform decisions.

\section{Future Goals}

The authors intend to carry forward the combination of deployable building and honest evaluation: extending the ROI construction (e.g.\ motion-compensated backgrounds, learned priors \cite{Zhao2025ASVC}), validating on broader scene sets, and pursuing the calibration and forensic extensions listed in Chapter~14.
